\section{Problem Definition}

We study a host with tenants \(T=\{1,\ldots,N\}\). Each tenant is an
independently configured \rocksdb instance with its own memtables, WAL, SST
files, flush path, compaction path, and foreground workload. Tenants do not
share logical key spaces, but their background work draws from shared host
resources. Time is divided into control windows \(w=1,2,\ldots\). In each
window, the controller assigns tenant \(i\) a background budget \(b_i(w)\). The
unit of this budget is harness-dependent: epoch runs change rate limits and
background-job settings between epochs, while continuous runs keep tenants alive
and change the rate-limiter budget through \ratecall.

The aggregate budget is fixed. If \(B\) is the host budget, a valid assignment
must satisfy \(\sum_i b_i(w)=B\), subject to the discrete budget tiers exposed by
the implementation. Without this conservation rule, a controller could improve
high-pressure tenants simply by spending more host capacity.

The controller is online. At the start of window \(w\), it may use only
observations from windows up to \(w-1\). The observation vector \(x_i(w-1)\)
can include demand signals such as offered write demand and completed write
throughput; pressure symptoms such as completion gap, write-tail latency,
pending-compaction proxies, and L0 pressure; background-work signals such as
compaction bytes and flush/compaction events; read-side measurements; and the
prior assignment. It may not include the next high-pressure set, future phase
id, or evaluator-only tier labels. SAKI also does not observe a single
ground-truth scalar called compaction debt; it observes symptoms of deferred
work and foreground pressure.

The objective is targeted relief for tenants currently under high write
pressure, measured by write-tail latency and write throughput, while preserving
attribution to online reallocation under a fixed budget. Because reallocation is
a tradeoff, collateral effects must also be visible: the evaluation reports
mid- and low-pressure effects, total throughput, and liveness failures. Each
decision therefore chooses \(b_i(w)\) from implementable budget levels using
\(x_i(w-1)\), conserving \(B\) under delayed feedback.

Several constraints sharpen the scope. Policies cannot read workload labels or
future drift schedules. Continuous runs actuate only the runtime rate limiter,
rather than every RocksDB option. Assignments occur at window boundaries, so
very fast drift may outrun a previous-window controller. SAKI is therefore not a
new compaction algorithm, an in-engine global scheduler, or an oracle
classifier. The definition leaves the ranking rule open; the next section
instantiates this online fixed-budget problem with a conserved rank-and-assign
framework and a transparent demand-aware ranking key.
